\chapter{Tumor Ablation}

\label{ch:ch81_tumor_ablation}

Tumour ablation is a group of minimally invasive techniques that destroy tumour tissue using extreme heat (radiofrequency or microwave), extreme cold (cryotherapy), or chemical injection. It is used as a curative or palliative treatment for selected tumours when surgery is not feasible or as an alternative to surgery.

---

\section{1. Patient Perspective: What Tumour Ablation Involves}

If you have a tumour that cannot be removed by surgery, your doctor may suggest \textbf{ablation} --- destroying the tumour with heat or cold using a needle-like probe.

\begin{itemize}
  \item The procedure is performed under local anaesthetic with sedation, or occasionally general anaesthesia.
  \item Under CT or ultrasound guidance, a thin probe is inserted through the skin directly into the tumour.
  \item The probe delivers energy (heat or cold) that destroys the tumour cells.
  \item The procedure typically takes 1--3 hours depending on tumour size and location.
  \item Most patients stay in hospital for 1--3 days. Pain and mild fever for a few days afterwards are common.
  \item Follow-up scans are arranged to confirm the tumour has been destroyed and to watch for recurrence.
\end{itemize}

---

\section{2. Ablation Techniques}

\begin{table}[h]
\centering
\begin{tabular}{l l l}
\hline
\textbf{Technique} & \textbf{Mechanism} & \textbf{Key Features} \\
\hline
Cryoablation & Argon gas freezes tissue to $-$40 to $-$200$^\circ$C, forming ice ball; thaw cycle destroys cells & Visible ice ball on CT/US; less pain; suitable near capsule surface \\
Microwave ablation (MWA) & Electromagnetic waves excite water molecules, generating heat to $>$ 150$^\circ$C & Faster; larger ablation zones; less affected by heat-sink effect \\
Radiofrequency ablation (RMA) & High-frequency alternating current heats tissue to 60--100$^\circ$C & Most widely used; predictable ablation zone; heat-sink effect near large vessels \\
\hline
\end{tabular}
\caption{Ablation Techniques}
\end{table}

\section{3. Indications by Organ}

\begin{table}[h]
\centering
\begin{tabular}{l l l}
\hline
\textbf{Organ} & \textbf{Indication} & \textbf{Typical Criteria} \\
\hline
Bone & Painful metastases, osteoid osteoma & Palliation of pain; local control \\
Kidney & Renal cell carcinoma & T1a ($\leq$ 4 cm), patient unfit for or refusing surgery \\
Liver & HCC, colorectal metastases, cholangiocarcinoma & $\leq$ 3 tumours, each $\leq$ 3 cm (Milan criteria for HCC);Child-Pugh A/B \\
Lung & Primary lung cancer or metastases & Peripheral lesion; $\leq$ 3 cm preferred; inoperable patient \\
\hline
\end{tabular}
\caption{Ablation Indications by Organ}
\end{table}

\begin{pearlbox}[Liver Ablation and the Heat-Sink Effect]
Tumours adjacent to large hepatic veins ($>$ 3 mm) or the IVC may have incomplete ablation due to the ``heat-sink effect'' --- blood flow carries away heat, reducing the temperature at the tumour margin. Microwave ablation is less susceptible to this phenomenon than RFA.
\end{pearlbox}

---

\section{4. Planning and Safety Considerations}

\begin{itemize}
  \item \textbf{Ablation margin}: Aim for a minimum 5 mm ablative margin beyond the visible tumour edge on follow-up imaging. A 1 cm margin is preferred for liver HCC.
  \item \textbf{Subcapsular location}: Tumours on the liver or kidney surface carry a higher risk of haemorrhage and incomplete ablation (no surrounding parenchyma to conduct heat). Consider artificial ascites or hydrodissection to displace adjacent organs.
  \item \textbf{Adjacent structures}: Map the proximity of bowel, gallbladder, diaphragm, chest wall, and nerves. Hydrodissection (injecting fluid between the tumour and adjacent structure) can create a safety margin.
  \item \textbf{Number of probes}: Multiple probes may be needed for larger tumours ($>$ 2 cm for RFA; $>$ 3 cm for MWA).
  \item \textbf{Anaesthesia}: Conscious sedation for most liver and kidney ablations; general anaesthesia for lung and bone ablations or complex cases.
\end{itemize}

\section{5. When to Recommend Ablation vs Resection}

\begin{itemize}
  \item \textbf{Ablation is preferred} when the patient is unfit for surgery, has limited hepatic reserve, has a small peripheral tumour, has multifocal disease not amenable to resection, or desires a minimally invasive approach.
  \item \textbf{Resection is preferred} for larger tumours ($>$ 3 cm for liver), tumours with vascular invasion, those in locations where ablation is technically challenging (e.g., central liver near major ducts), and when tissue diagnosis is required.
  \item \textbf{Combination}: Ablation may be used as a bridge to transplantation (e.g., liver HCC within Milan criteria) or combined with resection for multifocal disease.
\end{itemize}

---

\section{6. Follow-Up Imaging Protocol}

\begin{table}[h]
\centering
\begin{tabular}{l l l}
\hline
\textbf{Timepoint} & \textbf{Modality} & \textbf{What to Look For} \\
\hline
1 month post-ablation & Contrast-enhanced CT or MRI & Ablation zone (non-enhancing); assess for residual enhancing tumour at margin \\
3 months & Contrast-enhanced CT or MRI & Confirm complete necrosis; no nodular enhancement at margin \\
6 months then annually & Contrast-enhanced CT or MRI & Surveillance for local recurrence or new lesions \\
\hline
\end{tabular}
\caption{Post-Ablation Follow-Up Protocol (Liver)}
\end{table}

\begin{itemize}
  \item \textbf{Residual tumour}: Nodular enhancement at the ablation margin on the 1-month scan indicates incomplete treatment. Repeat ablation or consider alternative therapy.
  \item \textbf{Ablation zone evolution}: The ablation zone initially appears as a non-enhancing area with surrounding oedema. Over months, it contracts and fibroses.
  \item \textbf{Lung ablation follow-up}: Ground-glass opacity surrounding the ablation zone is expected and represents the ``halo sign'' of peri-ablational haemorrhage/inflammation. It should resolve over weeks.
\end{itemize}

\begin{pearlbox}[Clinical Pearl]
After liver RFA or MWA, the ablation zone on the 1-month scan should be \textit{larger} than the original tumour. If it is the same size or smaller, incomplete ablation should be suspected and correlation with arterial-phase enhancement is essential.
\end{pearlbox}

---

\section{7. Structured Reporting Example}

\begin{testbox}[Indications:]
68-year-old male with hepatitis C-related HCC. Single 2.2 cm lesion in segment VI of the liver. Child-Pugh A. MDT decision: percutaneous MWA as bridge to transplantation.
\end{testbox}
\begin{refbox}[Findings:]
\textbf{Pre-procedure}: Contrast-enhanced MRI shows a 2.2 cm arterially enhancing lesion in segment VI with washout on portal venous and delayed phases, consistent with HCC (LI-RADS 5). No vascular invasion. No other lesions.\\
\textbf{Procedure}: Under general anaesthesia and CT/ultrasound guidance, two microwave antenna probes were advanced into the lesion. Ablation performed at 100 W for 8 minutes per antenna. Intraprocedural ultrasound showed hyperechoic change throughout the tumour and 1 cm beyond.\\
\textbf{Post-procedure CT}: A 4.5 $\times$ 4.0 cm non-enhancing ablation zone in segment VI, encompassing the original tumour with a margin $>$ 5 mm. Small peri-ablational ground-glass halo. No haemorrhage, no pneumothorax, no biliary dilation.
\end{refbox}
\begin{pearlbox}[Impression:]
1. Successful microwave ablation of a 2.2 cm segment VI HCC.
2. Complete ablation achieved with an adequate ablative margin ($>$ 5 mm).
3. No immediate complications.
4. Follow-up contrast-enhanced MRI at 1 month to confirm complete response. Continue surveillance as bridge to transplant listing.
\end{pearlbox}

\clearpage
